\section{Introduction}
\label{sec:intro}

Large language model (LLM) serving is expensive for a reason that is easy to understate: most of the work is not ``running the model on the question,'' but \emph{writing the answer}. A typical chat request has two phases. \emph{Prefill} consumes the prompt in one (or a few) highly parallel forward passes of the serving Transformer: every prompt token is looked up in the model's embedding table and then transformed by every layer, on the accelerator. \emph{Decode} then emits tokens autoregressively; each step streams the same model weights---and, as the context grows, the key--value (KV) cache---from high-bandwidth memory. On contemporary accelerators this second phase is memory-bandwidth bound and has poor model-FLOP utilization~\citep{pope2023efficiently,leviathan2023fast}. Section~\ref{sec:cost} will show that a short chat prefill is often memory-bound as well: the difference is that prefill streams the weights \emph{once}, while decode streams them once \emph{per output token}. For the short-to-medium prompts and medium-to-long answers that dominate consumer chat, decode is the term that occupies the chip, sets time-to-last-token, and, once a fleet is provisioned for peak tokens-per-second, sets the bill.

That would be merely an unfortunate fact about Transformers if every request asked something new. In practice it does not. Search engines learned decades ago that query streams are heavy-tailed: a modest number of intents account for a large fraction of volume, each intent appearing as many surface forms. Public chat corpora show the same shape. Topic mass concentrates~\citep{zheng2024lmsys,zhao2024wildchat}; clusters of near-paraphrase are common; and for a surprising fraction of those clusters the \emph{acceptable answers} form a basin, not a point. ``How do I reverse a list in Python?'', ``python reverse list in place'', and ``what's the idiomatic way to flip a list?'' do not require three independently decoded completions. They require one good completion, reused.

This paper asks whether a provider that already sees this traffic---Google, or any other lab operating a frontier chatbot at scale---can stop paying decode for the head of that distribution. We call the resulting policy \emph{completion caching}. The object in the cache is not a KV block or a prompt prefix. It is a pair $(Q, C(Q))$: a previously seen query and the completion the model already produced for it. Given a new query $Q$, the system must decide whether some cached completion $C^\star$ is an \emph{acceptable} answer to $Q$. If so, it returns $C^\star$ and never decodes. A first decision can be made from a small retrieval embedder $e(Q)$, a nearest-neighbor lookup, and a gate---cheap enough to run on the host. If that gate is unsure, a serving-model prefill of $Q$ supplies a hidden state $h(Q)$ in the model's own geometry and a second gate can still skip decode. If that gate rejects, the model decodes as usual and the new pair may be inserted.

\subsection{What is already cached, and what is not}
\label{sec:intro-prior}

Serving stacks already cache aggressively. None of the standard caches is aimed at this decision.

Prefix and paged KV caches~\citep{kwon2023efficient} reuse attention state when two requests share a token prefix. Prompt-cache products do the same at the API boundary. These techniques are valuable: they cut prefill, improve time-to-first-token, and make multi-turn and templated workloads cheaper. They do not skip decode. After the shared prefix is restored, the model still writes a new answer, one bandwidth-bound step at a time. Prefill--decode disaggregation~\citep{zhong2024distserve} makes that fact architectural---the two phases even live on different machines---but it still \emph{pays} for decode.

Exact-match response caches (hash the prompt, return the bytes) are the other extreme. They skip both phases, but only when the user types the same string again, and they are unsafe even then: two users can type ``best Italian restaurant near me'' and must not receive each other's answer. Paraphrase without a shared context is the typical portable case; paraphrase \emph{with} a bound context is a different object, treated in Section~\ref{sec:method-bind}.

Semantic caches such as GPTCache~\citep{bang2023gptcache}, and more recent context-aware variants~\citep{yan2025contextcache,zhang2025smartcache}, are the closest prior work. They embed the query, retrieve a stored response by vector similarity, and, on a hit, bypass the model entirely. That is the right systems instinct, and production deployments have shown large latency wins \emph{when the hit is correct}. The limitation is the decision rule. Cosine similarity between query embeddings is a proxy for ``these questions are about the same thing.'' It is not a proxy for ``this stored completion is an acceptable answer to \emph{this} question.'' Near-neighbors in embedding space include near-misses that matter: different versions (``install Python~3.11'' versus ``3.12''), different entities with similar form (``capital of France'' versus ``capital of Spain''), and questions that share a topic but not an answer. Application-layer caches also typically treat the LLM as a black box. They do not use the model's own prefill state, they rarely degrade into a cheaper decode when the match is only partial, and they are seldom accompanied by a hardware cost model that says when a lookup is actually cheaper than generating.

A third line of work reuses \emph{text} without claiming the full answer is reusable. Retrieval-based speculative decoding~\citep{he2024rest} and prompt-lookup decoding draft tokens from a datastore or from the prompt itself; the model still verifies every token. These methods convert redundancy into a speedup with no quality change. They are complementary: a cached completion that is not safe to serve wholesale can still be a strong draft.

Completion caching sits between these lines. Like a semantic cache, it tries to skip decode. Like a serving-stack cache, it is willing to spend a prefill---a signal the model is already designed to compute---before making that call. Like speculative retrieval, it can fall back to using $C^\star$ as a draft rather than as an answer. The criterion it optimizes is not query similarity and not prefix equality. It is \emph{answer acceptability}.

\subsection{The idea}
\label{sec:intro-idea}

Write $C(\cdot)$ for the serving model's completion map, and write $\mathcal{A}(Q, C)$ for the (unobserved) predicate that $C$ is an acceptable reply to $Q$. A completion cache is a set of pairs $\{(Q_i, C_i)\}$ with $C_i = C(Q_i)$ at insertion time. On a new query $Q$ the system must either return some $C_i$ for which $\mathcal{A}(Q, C_i)$ holds, or report a miss and decode $C(Q)$.

$\mathcal{A}$ is not computable from $Q$ alone, and treating $\mathcal{A}(Q, C_i)$ as ``$Q$ is close to $Q_i$ in embedding space'' is exactly the error mode above. The useful observation is that we have richer cheap signals than that:
\begin{itemize}
  \item a retrieval embedding $e(Q)$ from a small dual encoder (tens to hundreds of millions of parameters), enough to retrieve a shortlist and, when the match is overwhelming, to return a cached completion without touching the serving model;
  \item the serving model's last-token prefill state $h(Q)$, already computed if decode is about to start, and a query representation in the model's own geometry;
  \item lightweight gates on $(e(Q), e(Q_i), h(Q), h(Q_i), C_i)$ that can be trained or calibrated to be conservative: a false hit serves the wrong answer, and is much more expensive than a false miss, which only costs a decode that would have been run anyway.
\end{itemize}
Those signals suggest a three-tier policy, specified in Section~\ref{sec:method}. A high-confidence, \emph{portable} hit never touches an accelerator. An ambiguous shortlist is resolved by prefilling $Q$ and gating on $h(Q)$; a pass still skips decode. A reject falls through to ordinary generation, optionally using the retrieved completion as a speculative draft. Context-bound intents (location, account, live search) are not served from another user's $C_i$; they may still skip decode by replaying a cached tool plan into a completion template.

Two further mechanisms are worth naming here because they change the hit-rate calculus. \emph{Top-K Projective Pooling} treats the cache not as a single nearest neighbor but as a small set of completions that can be selected among, voted on, or (in a more speculative design) fused. Many intents have several acceptable answers; requiring a unique nearest neighbor leaves hits on the table. \emph{Gating} is the opposite pressure: it exists to keep the served set inside $\mathcal{A}$. The systems problem here is transformed from maximizing cosine recall to maximizing decode-skips subject to a bound on unacceptable serves. \emph{Top-K Projective Pooling} is a technique for solving this problem; \emph{Gating} is a policy for choosing among the solutions.

\subsection{Why this is worth doing at provider scale}
\label{sec:intro-why}

Three facts have to be true at once for completion caching to matter. We argue that for a large consumer chatbot they are.

First, decode really is the scarce resource. Section~\ref{sec:cost} makes this quantitative across GPU and TPU parts, but the qualitative picture is stable. Prefill of a hundred-token question is one parallel pass. Decode of a three-hundred-token answer is hundreds of sequential, bandwidth-bound steps. Batching amortizes weight traffic across requests and is the main reason serving is viable; it does not make skipped decode free, because a skipped request is a slot that another request can occupy. At fleet scale, hit rate times decode cost is accelerator-hours.

Second, the lookup is cheap relative to that cost. Embedding a short query and searching an HNSW or IVF index of $10^7$--$10^9$ vectors is a CPU-and-DRAM problem, typically milliseconds, and is straightforward to colocate with the request router. Returning a few hundred tokens of cached text is a network problem measured in the same units. The comparison that matters is not ``is retrieval free?'' It is ``is retrieval plus a conservative gate cheaper than occupying an HBM-bound chip for the length of the answer?'' For any plausible chat length, yes---provided the hit is real.

Third, the traffic is skewed enough that a conservative gate still sees volume. We do not need, and should not claim, that ``most'' queries are duplicates. We need the head of the intent distribution---how-to questions, definitions, boilerplate code, exam-style problems, recipes, consumer product questions---to be large in absolute terms. A provider serving hundreds of thousands of queries per second can justify a substantial cache engineering effort at a single-digit-to-low-double-digit hit rate, if those hits are concentrated on decode-heavy, low-personalization requests. The long tail still decodes. That is expected; it is also how every successful cache works.

The same facts explain where the idea is \emph{not} worthwhile. Long-context workloads whose cost is prefill (summarize this document; answer from this unique corpus) gain little from a completion cache: the query is unique and decode may be the smaller term. Deeply personalized or tool-using turns (``what's on my calendar''; ``refactor \emph{this} file'') are not acceptable to serve from someone else's $C_i$. Time-sensitive questions go stale. A model-version bump invalidates the cache, or at least requires a freshness policy. These are design constraints, not objections to the head of the distribution.

\subsection{What this paper does}
\label{sec:intro-contrib}

This is a design paper. We do not yet report a production deployment. The contribution is to treat completion reuse as a first-class serving decision and to work out whether, and under what policy, it pays.

Section~\ref{sec:cost} develops a napkin cost model for prefill and decode on GPU and TPU, including the host and network terms a cache introduces. Section~\ref{sec:method} specifies the method: the serve path, gating, top-$k$ projective pooling, and the miss-side draft and insert rules. Section~\ref{sec:savings} translates hit rate and gate conservatism into fleet-level savings.

The claim we are willing to defend is modest and, we think, correct. At provider scale, a non-trivial fraction of chat decode is spent writing answers that already exist in the fleet's own history. Deciding that---carefully, mostly on CPU, and with prefill as a refinement rather than as a commitment to generate---is one of the few optimizations that can remove decode rather than merely accelerate it.
